% !TEX root = ../algebra_02.tex

\section{Степень расширения. Мультипликативность степени}

\begin{Definition}{Расширение поля}{}
	Пусть $K$ --- поле.
	$L$ называется \textbf{расширением} $K$, если $L$ --- поле и $K$ является подполем $L$. Обозначается как $L/K$.
\end{Definition}

Любое расширение $L$ поля $K$ можно рассматривать как линейное (векторное) пространство над $K$.
Операция умножения задана изначально как $L \times L \to L$, $(x,y) \mapsto x \cdot y$.
Ограничив её до $K \times L \to L$, получаем операцию умножения вектора из $L$ на скаляр из $K$.

\begin{Definition}{Конечное расширение}{}
	Расширение $L/K$ называется \textbf{конечным}, если $L$ конечномерно как векторное пространство над $K$.
	Размерность этого пространства $\dim_K L$ называется \textbf{степенью расширения} и обозначается $[L:K]$.
\end{Definition}

Расширение $L/K$ конечно тогда и только тогда, когда $[L:K] < \infty$.

\begin{Theorem}{О мультипликативности степени}{}
	Пусть даны конечные расширения $L/K$ и $K/F$. Тогда расширение $L/F$ также конечно, и выполняется равенство:
	\[ [L:F] = [L:K] \cdot [K:F] \]
\end{Theorem}

\begin{proof}
	Пусть $e_1, \dots, e_m$ --- базис $K$ над $F$, а $\delta_1, \dots, \delta_n$ --- базис $L$ над $K$. Докажем, что система элементов $\{ e_i \delta_j \mid i=1\dots m, j=1\dots n \}$ является базисом $L/F$.

	1. \textbf{Порождение:} Пусть $a \in L$. Так как $\delta_j$ образуют базис $L/K$, элемент $a$ можно разложить:
	\[ a = \sum_{j=1}^n b_j \delta_j, \quad b_j \in K \]
	Так как $e_i$ образуют базис $K/F$, каждый коэффициент $b_j$ можно разложить:
	\[ b_j = \sum_{i=1}^m c_{ij} e_i, \quad c_{ij} \in F \]
	Подставляя, получаем:
	\[ a = \sum_{j=1}^n \left( \sum_{i=1}^m c_{ij} e_i \right) \delta_j = \sum_{i,j} c_{ij} e_i \delta_j \]
	Таким образом, система $\{ e_i \delta_j \}$ порождает $L$ над $F$.

	2. \textbf{Линейная независимость:} Пусть некоторая линейная комбинация равна нулю:
	\[ \sum_{i,j} c_{ij} e_i \delta_j = 0, \quad c_{ij} \in F \]
	Сгруппируем слагаемые при $\delta_j$:
	\[ \sum_{j=1}^n \left( \sum_{i=1}^m c_{ij} e_i \right) \delta_j = 0 \]
	Выражение в скобках является элементом поля $K$. Так как элементы $\delta_j$ линейно независимы над $K$, все коэффициенты при них равны нулю:
	\[ \sum_{i=1}^m c_{ij} e_i = 0 \quad \text{для всех } j=1\dots n \]
	Так как элементы $e_i$ линейно независимы над $F$, отсюда следует, что:
	\[ c_{ij} = 0 \quad \text{для всех } i, j \]
	Следовательно, система линейно независима и образует базис.
\end{proof}
